\subsection{A partial short-map restriction certificate}

\begin{proposition}[Proved: filling the short two-cycles]
\label{A7:partial-short-restriction}
Fix a positive hom-set profile for an acyclic category on the full
chain $0<1<2<3<4$. Write
\[
 \begin{gathered}
 A=\mathcal C(0,1),\quad B=\mathcal C(1,2),\quad
 C=\mathcal C(2,3),\quad D=\mathcal C(3,4),\\
 U=\mathcal C(0,2),\quad V=\mathcal C(1,3),\quad
 W=\mathcal C(2,4).
 \end{gathered}
\]
Suppose the short map $P:A\times B\to U$ is completely assigned,
and an arbitrary subset of the values of $Q:B\times C\to V$ is
assigned. Work over $\mathbb F_2$. Define
\[
 d_{012}:\mathbb F_2[A\times B]\longrightarrow
 \mathbb F_2[A]\oplus\mathbb F_2[B]\oplus\mathbb F_2[U],
 \qquad (a,b)\longmapsto [a]+[b]+[P(a,b)].
\]
In the vector space with four distinguished summands
\[
 E=\mathbb F_2[A\times B]\oplus\mathbb F_2[A\times V]
       \oplus\mathbb F_2[U\times C]\oplus\mathbb F_2[B\times C],
\]
let $L$ be the span of the vectors
\[
 (a,b)+(a,Q(b,c))+(P(a,b),c)+(b,c)
 \tag{\(\dagger\)}
\]
for all $a\in A$ and all pairs $(b,c)$ whose $Q$ value is known.
If the embedded subspace $\ker d_{012}\subset E$ is contained
in $L$, then every associative completion of these short-map
assignments has zero products
\[
 H^2(B\mathcal C;\mathbb F_2)\otimes
 H^2(B\mathcal C;\mathbb F_2)\longrightarrow
 H^4(B\mathcal C;\mathbb F_2).
\]
In particular, no such completion models $\CP^2$.

There is a dual criterion: if $R:C\times D\to W$ is completely
assigned, and every two-cycle on the full subcategory $234$ is
a linear combination of the known $1234$ three-simplex
boundaries, then the same conclusion holds. The criterion does
not require a bound on any hom-set size, a top-homology
calculation, or an assignment of the longer compositions.
\end{proposition}
\begin{proof}
Consider any associative completion. The normalized nerve of its
full subcategory $012$ has no simplices of dimension three.
Consequently its second homology is exactly $\ker d_{012}$.
Each displayed vector \((\dagger)\) is the actual boundary of
the three-simplex $(a,b,c)$ in the completed category, with its
four faces lying respectively on $012,013,023,123$. Those faces
depend only on $P(a,b)$ and $Q(b,c)$; no longer composition is
used. Thus the assumed inclusion supplies a three-chain bounding
every two-cycle of $012$ in the completed nerve.

Let $\alpha$ be any global two-cocycle. It evaluates to zero on
every such boundary, hence on every two-cycle of $012$.
Over a field the annihilator of $\ker d_{012}$ is the image of
the transpose of $d_{012}$. Therefore
$\alpha|_{012}=\delta\lambda$ for a one-cochain $\lambda$ on
the edges of $012$. Extend $\lambda$ by zero to all other edges
and replace $\alpha$ by $\alpha-\delta\lambda$. This new global
representative vanishes on every two-simplex on $012$.

Every nondegenerate four-simplex of an acyclic five-object
category has object string $01234$. For any other two-cocycle
$\beta$, its Alexander--Whitney product with the new
representative of $\alpha$ is consequently zero on every
four-simplex. This proves the asserted vanishing of all
degree-two products. In the dual argument one changes the
right factor to vanish on $234$.

The mod-two cohomology of $\CP^2$ has a degree-two class with
nonzero square, so such a category cannot model $\CP^2$.
\end{proof}

\begin{remark*}
The partial data are used only to display vectors which become
actual three-boundaries in every associative completion. No
claim that the unspecified longer compositions already form a
category is needed. A certificate gives a complete basis of
$\ker d_{012}$ and, for each basis vector, an explicit linear
combination of the known three-boundaries equal to it in $E$.
The checker independently verifies the kernel dimension, basis
independence, every two-cycle equation, every filling equation,
and the absence of coefficients on unknown three-simplices.
An empty basis is allowed precisely when the checked kernel is
zero. All literal completions of a certified global prefix are
excluded, with the same unrestricted suffix-count convention as
the earlier prefix certificates.
\end{remark*}
